\section{模形式}\label{sec:modular-1-2}
$\Gamma$ 始终表示 $\mathrm{SL}_2(\mathbb Z)$ 的有限指数子群.
\subsection{权为零的模函数}\label{sec:modular-1-2-1}
回顾，Riemann 曲面 $X$ 上的亚纯函数 $X\dashrightarrow\mathbb C$（定义在其极点以外）唯一延拓为全纯映射 $X\longrightarrow\mathbb P^1(\mathbb C)$.记 $\mathcal M_X$ 为 $X$ 上亚纯函数构成的域.例如 $\mathcal M_{\mathbb P^1(\mathbb C)}\simeq\mathbb C(z)$.这一理论中有一个引人注目的结果（我们并不需要用到）：$X\mapsto\mathcal M_X$ 给出紧 Riemann 曲面及其间非常值全纯映射构成的范畴与 $\mathbb C$ 上有限生成、超越次数为 $1$ 的域构成的范畴之间的等价.\footnote{译注：按函数的拉回，这里通常应表述为反等价；原文称为等价.}这提供了研究域 $\mathcal M_{X(\Gamma)}$ 的动机.

设 $\varphi:X(\Gamma)\longrightarrow\mathbb P^1(\mathbb C)$ 为亚纯函数，记复合映射 $\mathbb H\longrightarrow X(\Gamma)\xrightarrow{\varphi}\mathbb P^1(\mathbb C)$ 为 $f$.这是 $\mathbb H$ 上 $\Gamma$ 不变的亚纯函数.反过来，给定这样的 $f$，它可以下降到 $Y(\Gamma)$，但要延拓到 $X(\Gamma)$，还需要在尖点处亚纯的条件.

为明确这一条件，先考虑较一般的情形：$\mathbb H$ 上的亚纯函数 $f$ 具有水平方向的周期，即存在非零整数 $h$，使 $f(z)=f(z+h)$（换言之，$f$ 在矩阵 $\left[\begin{smallmatrix}1&h\\0&1\end{smallmatrix}\right]$ 下不变）.它可写成 $q=\exp(2i\pi z/h)\in\mathbb D\setminus\{0\}$ 的亚纯函数，因而在 $q=0$（即 $z=\infty$）附近具有展开
\[
f(z)=\sum_{n\in\mathbb Z}a_nq^n.
\]
若这个关于 $q$ 的函数在 $0$ 处亚纯，即当 $n\ll0$ 时 $a_n=0$，则称 $f$ 在无穷远处亚纯.此时记
\[
\operatorname{ord}_{\infty,h}(f)=\text{使 }a_n\ne0\text{ 的最小整数 }n.
\]
若 $\operatorname{ord}_{\infty,h}(f)\ge0$，则称 $f$ 在无穷远处全纯.这些概念不依赖于使 $f$ 在 $\left[\begin{smallmatrix}1&h\\0&1\end{smallmatrix}\right]$ 下不变的 $h$ 的选择（当然整数 $\operatorname{ord}_{\infty,h}$ 依赖于 $h$）.

再来看 $\Gamma$ 不变的亚纯函数.对尖点 $x\in\mathbb P^1(\mathbb Q)$，选取 $\tau\in\mathrm{SL}_2(\mathbb Z)$ 使 $\tau x=\infty$；若具有水平周期的函数 $\tau^*f:z\mapsto f(\tau^{-1}z)$ 在无穷远处亚纯或全纯，则相应称 $f$ 在尖点 $x$ 处亚纯或全纯.这不依赖于 $\tau$ 的选择，整数
\[
\operatorname{ord}_{x,h}(f):=\operatorname{ord}_{\infty,h}(\tau^*f)
\]
也一样.$\mathbb H$ 上 $\Gamma$ 不变且在 $\mathbb H$ 上及各尖点处亚纯的函数 $f:\mathbb H\longrightarrow\mathbb C$ 称为级为 $\Gamma$、权为 $0$ 的\newterm{modular-modular-function}{模函数}{modular function}.它们对应于 $X(\Gamma)$ 上的亚纯函数.

$\mathbb H$ 上 $\Gamma$ 不变且在 $\mathbb H$ 上及各尖点处全纯的函数称为级为 $\Gamma$、权为 $0$ 的\newterm{modular-modular-form}{模形式}{modular form}.它们对应于 $X(\Gamma)$ 上的全纯函数，因此……都是常值函数！
\subsection{权为二的模形式}\label{sec:modular-1-2-2}
在亚纯函数之后，自然应考察 $X(\Gamma)$ 上的亚纯微分形式.在 $\mathbb C$ 的开集 $U$ 上，全纯（分别地，亚纯）微分形式写作 $\omega=f(z)dz$，其中 $f$ 全纯（分别地，亚纯）.若 $\varphi:V\xrightarrow{\sim}U$ 是从 $\mathbb C$ 的另一开集出发的双全纯映射，则可将 $\omega$ 拉回为 $\varphi^*\omega=f(\varphi(z))\varphi'(z)dz$.在一般的 Riemann 曲面 $X$ 上，一个微分形式是在每张坐标图 $X\supseteq O\xrightarrow{\sim}U\subset\mathbb C$ 中给定一个 $U$ 上的微分形式，并要求它们与上述坐标变换公式相容.

设 $\pi:\mathbb H\longrightarrow X(\Gamma)$ 是典范投影在 $\mathbb H\subset\mathbb H^*$ 上的限制.$X(\Gamma)$ 上的亚纯微分形式 $\omega$ 诱导 $\Gamma$ 不变的亚纯微分形式 $\pi^*\omega=f(z)dz$.不变性条件为
\[
\forall\gamma\in\Gamma,\qquad f(\gamma z)d(\gamma z)=f(z)dz.
\]
对 $\gamma=\left[\begin{smallmatrix}a&b\\c&d\end{smallmatrix}\right]$，令 $j_\gamma(z)=cz+d$，计算可得
\[
d(\gamma z)=j_\gamma(z)^{-2}dz.
\]
因此函数 $f$ 满足如下``自守性''：
\[
\forall\gamma\in\Gamma,\qquad f(z)=j_\gamma(z)^{-2}f(\gamma z).
\]
由于 $j_{\pm\left[\begin{smallmatrix}1&h\\0&1\end{smallmatrix}\right]}(z)^2=1$，每个满足这一自守性的函数 $f$ 都具有水平周期.同样，对任意 $\tau\in\mathrm{SL}_2(\mathbb Z)$，微分形式 $\tau^*(f\,dz)=j_\tau(z)^{-2}f(\tau z)dz$ 在 $\tau^{-1}\Gamma\tau$ 下不变，故函数 $j_\tau(z)^{-2}f(\tau z)$ 具有水平周期.若此函数在无穷远处亚纯或全纯，就相应称 $f$ 在尖点 $\tau\infty$ 处亚纯或全纯.这不依赖于 $\tau$ 的选择，整数\footnote{原注：应注意，这个 $\operatorname{ord}_{x,h}(f)$ 的定义不同于先前针对 $\Gamma$ 不变函数的定义（除非 $x=\infty$）.这里是 $f\,dz$ 在 $\Gamma$ 下不变.记号似乎有歧义，但只有当 $f$ 与 $f\,dz$ 同时 $\Gamma$ 不变时才会有歧义，而这样的函数只能是零.}
\[
\operatorname{ord}_{x,h}(f):=\operatorname{ord}_{\infty,h}\bigl(j_\tau(z)^{-2}f(\tau z)\bigr)
\]
也一样.

级为 $\Gamma$、权为 $2$ 的模函数是满足上述自守性，且在 $\mathbb H$ 上和各尖点处亚纯的函数.因此它对应于 $X(\Gamma)$ 上的亚纯微分形式.级为 $\Gamma$、权为 $2$ 的模形式是进一步在 $\mathbb H$ 上和各尖点处全纯的模函数.

\textbf{注意.}$f$ 的全纯性与 $\omega$ 的全纯性并不等价！为解释这一点，记 Riemann 曲面上亚纯函数 $f$ 在点 $P$ 处的零点或极点阶数为 $\operatorname{ord}_P(f)$.若在 $P$ 附近的局部坐标 $z$ 中写成 $f=g(z)$，则 $\operatorname{ord}_P(f)=\operatorname{ord}_0(g)$.同样，若亚纯微分形式写成 $\omega=g(z)dz$，则整数 $\operatorname{ord}_0(g)$ 不依赖于坐标选择，记为 $\operatorname{ord}_P(\omega)$.当然，$\omega$ 在 $P$ 处全纯当且仅当 $\operatorname{ord}_P(\omega)\ge0$.
\begin{lemma*}
设 $\omega$ 是 $X(\Gamma)$ 上的亚纯微分形式，$f$ 是 $\mathbb H$ 上满足 $\pi^*\omega=f(z)dz$ 的亚纯函数，则：
\begin{enumerate}[label=\roman*)]
\item 若 $P\in\mathbb H$ 是 $\Gamma$ 的普通点，则 $\operatorname{ord}_P(f)=\operatorname{ord}_{\pi(P)}(\omega)$.
\item 若 $P$ 是阶为 $e$（等于 $2$ 或 $3$）的椭圆点，则 $\operatorname{ord}_P(f)=e\operatorname{ord}_{\pi(P)}(\omega)+e-1$.
\item 若 $x$ 是阶为 $h$ 的尖点，则 $\operatorname{ord}_{x,h}(f)=\operatorname{ord}_{\pi(x)}(\omega)+1$.
\end{enumerate}
\end{lemma*}
\begin{proof}
(i)、(ii)：前面已经看到，在 $\mathbb H$ 的点 $P$ 附近存在局部坐标 $z'=\tau(z)$，使 $\pi$ 具有 $z'\mapsto z'^e$ 的形式，普通点处 $e=1$.因此 $u=z'^e$ 下降为 $\pi(P)$ 附近的局部坐标.若在 $\pi(z_0)$ 附近 $\omega=g(u)du$，则 $\pi^*\omega=g(z'^e)d(z'^e)=ez'^{e-1}g(z'^e)dz'$，故
\begin{align*}
\operatorname{ord}_P(f)&=\operatorname{ord}_P(\pi^*\omega)=\operatorname{ord}_0(z'^{e-1}g(z'^e))\\
&=e\operatorname{ord}_0(g)+e-1=e\operatorname{ord}_{\pi(P)}(\omega)+e-1.
\end{align*}
(iii)：必要时用 $\tau\in\mathrm{SL}_2(\mathbb Z)$ 作变换，可设 $x=\infty$.这时 $q=\exp(2i\pi z/h)$ 是 $\pi(x)$ 附近的局部坐标.若 $\omega=g(q)dq$，则
\[
\pi^*\omega=\frac{2i\pi}{h}g(\exp(2i\pi z/h))\exp(2i\pi z/h)dz=f(z)dz.
\]
于是 $f(z)=\frac{2i\pi}{h}qg(q)$，最终
\[
\operatorname{ord}_{\infty,h}(f)=\operatorname{ord}_0(qg(q))=\operatorname{ord}_0(g(q))+1=\operatorname{ord}_{\pi(x)}(\omega)+1.
\]
\end{proof}
记 $\mathcal M_2(\Gamma)$ 为级为 $\Gamma$、权为 $2$ 的模形式构成的 $\mathbb C$ 向量空间；记 $\mathcal S_2(\Gamma)$ 为其中\newterm{modular-cusp-form}{尖点形式}{cusp form}构成的子空间（法语也称 formes paraboliques 或 cuspidales）.按定义，它们在所有尖点处为零，即每个尖点 $x$ 都满足 $\operatorname{ord}_{x,h}(f)>0$.

Weil 除子理论，特别是 Riemann--Roch 定理，表明这些空间有限维，并能计算其维数.先回顾一些事实.
\begin{itemize}
\item $X$ 上的除子是和 $\sum_{P\in X}a_P[P]$，其中 $a_P\in\mathbb Z$，只有有限多个系数非零.因此除子全体 $\operatorname{Div}(X)$ 是以 $X$ 为基的自由 Abel 群.它带有偏序：$\sum_Pa_P[P]\le\sum_Pb_P[P]$ 当且仅当所有 $P$ 都满足 $a_P\le b_P$.
\item 除子的次数定义为 $\deg(\sum_Pa_P[P])=\sum_Pa_P\in\mathbb Z$.
\item 每个非零亚纯函数 $f$ 都对应除子
\[
\operatorname{div}(f):=\sum_P\operatorname{ord}_P(f)[P].
\]
有 $\operatorname{div}(fg)=\operatorname{div}(f)+\operatorname{div}(g)$ 和 $\operatorname{div}(f+g)\ge\inf(\operatorname{div}(f),\operatorname{div}(g))$.\footnote{原注：这里的 $\inf$ 应理解为除子偏序的``最大公因子''，即同时小于等于 $\operatorname{div}(f)$ 和 $\operatorname{div}(g)$ 的元素中的最大者.}由此得到群同态 $\mathcal M_X^\times\xrightarrow{\operatorname{div}}\operatorname{Div}(X)$，其像记为 $\operatorname{Div.pr}(X)$，像中的除子称为主除子.商群 $\operatorname{Div}(X)/\operatorname{Div.pr}(X)$ 称为 $X$ 的 Picard 群，记为 $\operatorname{Pic}(X)$（它类似于数域的理想类群）.
\item 映射 $\deg$ 通过 $\operatorname{Pic}(X)$ 分解.事实上，将 $f\in\mathcal M_X^\times$ 看作全纯映射 $X\longrightarrow\mathbb P^1(\mathbb C)$，记其在 $P\in X$ 处的分歧指数为 $e_P$，则
\[
\operatorname{div}(f)=\sum_{P\mapsto0}e_P[P]-\sum_{P\mapsto\infty}e_P[P],
\]
故其次数为零.
\item 类似地，每个亚纯微分形式对应除子
\[
\operatorname{div}(\omega)=\sum_P\operatorname{ord}_P(\omega)[P].
\]
由于对任意 $f\in\mathcal M_X^\times$ 都有 $\operatorname{div}(f\omega)=\operatorname{div}(f)+\operatorname{div}(\omega)$，$\omega$ 在 $\operatorname{Pic}(X)$ 中的像不依赖于 $\omega$，称为 $X$ 的典范除子，记为 $K$.
\item 若 $D\in\operatorname{Div}(X)$，记
\[
L(D):=\{f\in\mathcal M_X^\times:\operatorname{div}(f)+D\ge0\}\sqcup\{0\}.
\]
这是 $\mathbb C$ 向量空间，其维数 $\ell(D):=\dim_{\mathbb C}L(D)$ 只依赖于 $D$ 在 $\operatorname{Pic}(X)$ 中的像，因为对 $f_0\in\mathcal M_X^\times$，乘以 $f_0$ 诱导线性同构 $L(D)\xrightarrow{\sim}L(D+\operatorname{div}(f_0))$.\footnote{译注：原文同构及其所称``乘以 $f_0$''均照录；按 $L(D)$ 的定义，这一方向应乘以 $f_0^{-1}$.}
\item Riemann--Roch 定理断言
\[
\ell(D)=\deg(D)+1-g+\ell(K-D),
\]
其中 $g$ 是 $X$ 的亏格.由 $\ell(0)=1$（紧曲面 $X$ 上的全纯函数均为常值）分别取 $D=0$ 和 $D=K$，特别得到
\[
\ell(K)=g,\qquad\deg(K)=2g-2.
\]
因此，若 $\deg(D)>2g-2$，则 $\deg(K-D)<0$，故 $L(K-D)=\{0\}$（因为有效除子的次数非负），从而
\[
\ell(D)=\deg(D)+1-g.
\]
\end{itemize}
回到 $X=X(\Gamma)$，记
\[
D_\pi:=\sum_{P\in X(\Gamma)}(1-1/e_P)[P]\in\operatorname{Div}(X(\Gamma))\otimes_{\mathbb Z}\mathbb Q,
\]
其中普通点处 $e_P=1$，椭圆点处 $e_P=2$ 或 $3$，尖点处 $e_P=+\infty$.于是 $\deg(D_\pi)=\frac12n_2+\frac23n_3+n_\infty$.
\begin{corollary*}
有等式
\[
\dim_{\mathbb C}\mathcal M_2(\Gamma)=\ell(D_\pi+K)=g-1+n_\infty,\qquad
\dim_{\mathbb C}\mathcal S_2(\Gamma)=\ell(K)=g.
\]
\end{corollary*}
\begin{proof}
前一引理表明，由 $\pi^*\omega=f(z)dz$ 定义的映射 $\omega\mapsto f$ 诱导同构
\[
\{\omega:\operatorname{div}(\omega)+D_\pi\ge0\}\sqcup\{0\}\xrightarrow{\sim}\mathcal M_2(\Gamma).
\]
固定 $X(\Gamma)$ 上的一个亚纯微分形式 $\omega_0$，\footnote{原注：这样的形式总是存在.例如在 $X(1)=\mathbb P^1$ 上可取 $dz$（其除子为 $\operatorname{div}(dz)=-2[\infty]$），然后沿 $X(\Gamma)\longrightarrow X(1)$ 拉回.}将其余形式写成 $\omega=g\omega_0$.由于 $\operatorname{div}(\omega)=\operatorname{div}(g)+\operatorname{div}(\omega_0)$，映射 $g\mapsto g\omega_0\mapsto f$ 诱导同构
\[
L(\operatorname{div}(\omega_0)+D_\pi)\xrightarrow{\sim}\mathcal M_2(\Gamma).
\]
同一映射还诱导同构
\[
L\Bigl(\operatorname{div}(\omega_0)+D_\pi-\sum_{P\mapsto\infty}[P]\Bigr)\xrightarrow{\sim}\mathcal S_2(\Gamma).
\]
记 $D_\pi$ 的整部为 $\lfloor D_\pi\rfloor$，即 $\lfloor D_\pi\rfloor=\sum_{P\mapsto\infty}[P]$.显然 $L(\operatorname{div}(\omega_0)+D_\pi)=L(\operatorname{div}(\omega_0)+\lfloor D_\pi\rfloor)$，于是
\begin{align*}
\dim_{\mathbb C}\mathcal M_2(\Gamma)&=\ell(\operatorname{div}(\omega_0)+\lfloor D_\pi\rfloor)\\
&=\ell(K+\lfloor D_\pi\rfloor)=\ell\Bigl(K+\sum_{P\mapsto\infty}[P]\Bigr).
\end{align*}
由于 $\deg(K+\sum_{P\mapsto\infty}[P])=2g-2+n_\infty>2g-2$，可用 Riemann--Roch 定理的简化形式得到
\[
\ell\Bigl(K+\sum_{P\mapsto\infty}[P]\Bigr)=\deg\Bigl(K+\sum_{P\mapsto\infty}[P]\Bigr)+1-g=g-1+n_\infty.
\]
同理得到
\[
\dim_{\mathbb C}\mathcal S_2(\Gamma)=\ell\Bigl(\operatorname{div}(\omega_0)+\lfloor D_\pi\rfloor-\sum_{P\mapsto\infty}[P]\Bigr)=\ell(K)=g.
\]
\end{proof}
\subsection{权为自然数 $k$ 的模形式}\label{sec:modular-1-2-3}
对整数 $k\in\mathbb Z$，用公式
\[
f[\gamma]_k(z):=\det(\gamma)^{k/2}j_\gamma(z)^{-k}f(\gamma z)
\]
定义 $\mathrm{GL}_2(\mathbb R)^+$ 在 $\mathcal M_X$ 上的作用.\footnote{译注：原文写作 $\mathcal M_X$；这里的公式实际针对 $\mathbb H$ 上的亚纯函数.}这确为作用，即 $[\gamma\gamma']_k=[\gamma]_k[\gamma']_k$，由容易验证的公式 $j_{\gamma\gamma'}(z)=j_\gamma(\gamma'z)j_{\gamma'}(z)$ 得出.希望上一小节已使下面的定义显得更自然.
\begin{definition*}
设 $\Gamma$ 为 $\mathrm{SL}_2(\mathbb Z)$ 的有限指数子群.级为 $\Gamma$、权为 $k\in\mathbb N$ 的模形式是满足以下条件的函数 $f:\mathbb H\longrightarrow\mathbb C$：
\begin{enumerate}[label=\roman*)]
\item $\forall\gamma\in\Gamma,\ f[\gamma]_k=f$（自守性）.
\item $f$ 在 $\mathbb H$ 上及各尖点处全纯.\footnote{原注：与权为 $2$ 时一样，由于 $f$ 具有水平周期，在尖点 $\infty$ 处全纯有明确定义.对尖点 $x=\tau\infty$，$\tau\in\mathrm{SL}_2(\mathbb Z)$，若同样具有水平周期的变换 $f[\tau]_k$ 在无穷远处全纯或亚纯，则相应称 $f$ 在 $x$ 处全纯或亚纯.这不依赖于 $\tau$ 的选择，整数 $\operatorname{ord}_{x,h}(f):=\operatorname{ord}_{\infty,h}(f[\tau]_k)$ 也一样.}
\end{enumerate}
这些函数构成 $\mathbb C$ 向量空间，记为 $\mathcal M_k(\Gamma)$.若 $f$ 在各尖点处为零，则称其为尖点形式，其空间记为 $\mathcal S_k(\Gamma)$.
\end{definition*}
若将这些条件中的``全纯''改为``亚纯''，则称满足条件的函数为级为 $\Gamma$、权为 $k$ 的模函数.
\begin{remark*}
若 $k$ 为奇数且 $-1\in\Gamma$，则 $\mathcal M_k(\Gamma)=\mathcal S_k(\Gamma)=0$，因为 $f[-1]_k=-f$.
\end{remark*}
\begin{remark*}
验证函数对于 $(\Gamma,k)$ 自守，只需对 $\Gamma$ 的一组生成元验证 $f[\gamma]_k=f$.例如 $\Gamma=\Gamma(1)$ 且 $k$ 为偶数时，只需验证 $f(z+1)=f(z)$ 与 $f(-1/z)=z^kf(z)$.
\end{remark*}
\begin{remark*}
若 $f$ 亚纯且为权为 $k$ 的自守函数，验证它在各尖点处全纯，只需证明对任意 $\tau\in\Gamma(1)$，函数 $\tau^*f(z)$ 在 $\Im(z)\to\infty$ 时存在极限.\footnote{译注：原文在这里使用 $\tau^*f$；结合上面的定义，它指带权变换 $f[\tau]_k$.}
\end{remark*}
\begin{remark*}
函数乘法诱导乘积 $\mathcal M_k(\Gamma)\otimes\mathcal M_{k'}(\Gamma)\longrightarrow\mathcal M_{k+k'}(\Gamma)$，使 $\mathcal M(\Gamma):=\bigoplus_{k\in\mathbb N}\mathcal M_k(\Gamma)$ 成为分次 $\mathbb C$ 代数，而 $\mathcal S(\Gamma):=\bigoplus_{k\in\mathbb N}\mathcal S_k(\Gamma)$ 为其中的理想.
\end{remark*}
\begin{remark*}
当 $\Gamma=\Gamma(1)$ 时，可用格上的函数表述自守性.$\mathcal R$ 上的函数称为权为 $k$ 的齐次函数，若
\[
\forall\Lambda\in\mathcal R,\quad\forall\alpha\in\mathbb C^\times,\quad\varphi(\alpha\Lambda)=\alpha^{-k}\varphi(\Lambda).
\]
记 $\Lambda_z:=z\mathbb Z\oplus\mathbb Z$.等式 $\Lambda_{\gamma z}=j_\gamma(z)^{-1}\Lambda_z$ 表明，由 $f(z)=\varphi(\Lambda_z)$ 定义的映射 $\varphi\mapsto f$ 给出双射
\[
\{\mathcal R\text{ 上权为 }k\text{ 的齐次函数}\}\longleftrightarrow\{\text{权为 }k\text{ 的 }\Gamma(1)\text{ 自守函数}\}.
\]
其逆映射为 $\varphi(\Lambda)=\omega_2^{-k}f(\omega_1/\omega_2)$，其中 $\binom{\omega_1}{\omega_2}$ 是 $\Lambda$ 的任意一组满足 $\omega_1/\omega_2\in\mathbb H$ 的基.
\end{remark*}
\begin{example*}[{Eisenstein 级数}]
构造 $\mathcal R$ 上权为 $k$ 的函数，最直接的方法是令
\[
G_k(\Lambda):=\sum_{\omega\in\Lambda^*}\frac1{\omega^k},\qquad\Lambda^*=\Lambda\setminus\{0\}.
\]
当 $k\ge3$ 时此和绝对收敛（习题），且如所应有，$k$ 为奇数时它为零.对应的 $\mathbb H$ 上的函数为
\[
G_k(z)=\sum_{(m,n)\in(\mathbb Z^2)^*}\frac1{(mz+n)^k}.
\]
为研究全纯性，取充分小的 $\varepsilon>0$，使存在 $z_\varepsilon$ 满足 $|z_\varepsilon|=1-\varepsilon$ 与 $2\Re(z_\varepsilon)=1+\varepsilon$.考虑基本域 $\overline D$ 的开邻域 $D^\varepsilon$，由不等式 $|z|>1-\varepsilon$ 与 $|2\Re(z)|<1+\varepsilon$ 给出.对 $z\in D^\varepsilon$，不等式
\[
|mz+n|^2=m^2z\bar z+2mn\Re(z)+n^2\ge m^2(1-\varepsilon)^2-mn(1+\varepsilon)+n^2=|mz_\varepsilon-n|^2
\]
表明 $G_k$ 在 $D^\varepsilon$ 上正规收敛，因而在那里全纯.\footnote{译注：原文不等式照录.若 $mn<0$，其中 $-mn(1+\varepsilon)$ 不是统一下界；通常应以 $-|mn|(1+\varepsilon)$ 作估计，并按符号选取比较式.}由自守性，$G_k$ 在整个 $\mathbb H$ 上全纯.要证明它在尖点 $\infty$ 处也全纯，只需证明 $\Im(z)\to\infty$ 时 $G_k(z)$ 存在极限.由 $T$ 不变性，可让 $z$ 留在 $\overline D$ 中；正规收敛性表明，由于 $k$ 为偶数，
\[
\lim_{z\in\overline D,\,\Im(z)\to\infty}G_k(z)=\sum_{n\in\mathbb Z\setminus\{0\}}\frac1{n^k}=2\zeta(k).
\]
因此 $G_k\in\mathcal M_k(\Gamma(1))$，且 $G_k(\infty)=2\zeta(k)$（仍假设 $k$ 为偶数）.这些称为\newterm{modular-eisenstein-series}{Eisenstein 级数}{Eisenstein series}.
\end{example*}
\begin{example*}[{$\Delta$ 函数}]
由于我们讨论椭圆曲线时才会说明的原因，习惯上令 $g_2:=60G_4$ 和 $g_3:=140G_6$.暂时承认 $\zeta(4)=\frac{\pi^4}{2\cdot3^2\cdot5}$ 与 $\zeta(6)=\frac{\pi^6}{3^3\cdot5\cdot7}$（见习题单），便可看出
\[
\Delta:=g_2^3-27g_3^2\in\mathcal S_{12}(\Gamma(1)),
\]
这给出了尖点形式的第一个例子.
\end{example*}
\subsection{模函数的极点与零点}\label{sec:modular-1-2-4}
当 $k=2k'$ 为偶数时，可将权为 $k$ 的模形式解释为 $X(\Gamma)$ 上某些 $k'$ 次微分形式.$\mathbb C$ 的开集上的 $k'$ 次微分形式是形如 $\omega=f(z)(dz)^{k'}$ 的表达式.若 $\varphi$ 是两个开集之间的双全纯映射，则通过公式 $\varphi^*\omega=f(\varphi(z))\varphi'(z)^{k'}(dz)^{k'}$ 拉回.由此可在流形上定义 $k'$ 次微分形式：在 $X$ 的每张坐标图上给定这样的形式，并与坐标变换相容.于是与权为 $2$ 时类似，特别有对应
\[
\{\text{级为 }\Gamma\text{、权为 }k\text{ 的模函数 }f\}\longleftrightarrow
\{X(\Gamma)\text{ 上亚纯的 }k'\text{ 次微分形式 }\omega\},
\]
由等式 $\pi^*\omega=f(z)(dz)^{k'}$ 给出.与权为 $2$ 时相同，$f$ 与 $\omega$ 的零点和极点满足
\[
\begin{cases}
\operatorname{ord}_P(f)=e\operatorname{ord}_{\pi(P)}(\omega)+k'(e-1),&P\in\mathbb H,\ e\text{ 为其分歧指数},\\
\operatorname{ord}_{x,h}(f)=\operatorname{ord}_{\pi(x)}(\omega)+k',&x\text{ 为阶为 }h\text{ 的尖点}.
\end{cases}
\]
\begin{proposition*}
设 $f$ 是级为 $\Gamma$、权为 $k$ 的非零模函数，则
\begin{align*}
\sum_{P\in\Gamma\backslash\mathbb H}\frac{\operatorname{ord}_P(f)}{e_P}+\sum_{x\in\Gamma\backslash P_\Gamma}\operatorname{ord}_{x,h_x}(f)
&=\frac k2\Bigl(2g-2+\frac12n_2+\frac23n_3+n_\infty\Bigr)\\
&=\frac{kd}{12},
\end{align*}
其中回顾 $d=[\Gamma(1):\Gamma\{\pm1\}]$.
\end{proposition*}
\begin{proof}
若 $k=2k'$ 为偶数，上述等式表明左端等于 $\deg(\operatorname{div}(\omega))+k'\deg(D_\pi)$.次数 $\deg(\operatorname{div}(\omega))$ 不依赖于 $\omega$（因为其余形式均为 $g\omega$，$g\in\mathcal M_{X(\Gamma)}^\times$）.计算时可取形如 $\omega'=\omega_1^{k'}$ 的形式，其中 $\omega_1$ 是 $X(\Gamma)$ 上的亚纯一次微分形式；$\omega_1^{k'}$ 指在 $\omega_1=f(z)dz$ 的坐标图上由 $\omega_1^{k'}=f(z)^{k'}(dz)^{k'}$ 给出的 $k'$ 次微分形式.得到 $\deg(\operatorname{div}(\omega))=k'\deg(\operatorname{div}(\omega_1))=k'(2g-2)$，由此推出第一个等式.第二个等式来自亏格公式.若 $k$ 为奇数，只需对 $f^2$ 应用偶数情形.
\end{proof}
\paragraph{应用：计算 $\mathcal M(\Gamma(1))$.}
对 $\Gamma=\Gamma(1)$，公式成为
\[
\operatorname{ord}_{\infty,1}(f)+\frac12\operatorname{ord}_i(f)+\frac13\operatorname{ord}_j(f)+\sum_{P\ne i,j,\infty}\operatorname{ord}_P(f)=k/12.
\]
已知 $k$ 为奇数时 $\mathcal M_k(\Gamma(1))$ 为零.取 $k=2$，由于每个 $\operatorname{ord}_{?}(f)\ge0$，可见不存在满足此等式的 $f$，故
\[
\mathcal M_2(\Gamma(1))=0.
\]
取 $k=4$，得到每个 $f\in\mathcal M_4(\Gamma(1))$ 在 $j$ 处有一个单零点，在其他地方不为零.这也适用于 $G_4$，并表明 $f/G_4$ 是权为 $0$ 的模形式，因而为常值，故
\[
\mathcal M_4(\Gamma(1))=\mathbb C G_4.
\]
同样取 $k=6$，得到 $G_6$ 在 $i$ 处有一个单零点，在其余地方不为零，从而
\[
\mathcal M_6(\Gamma(1))=\mathbb C G_6.
\]
同样的推理表明 $\mathcal M_8$ 与 $\mathcal M_{10}$ 均为一维；更确切地，
\[
\mathcal M_8(\Gamma(1))=\mathbb C G_8=\mathbb C G_4^2,\qquad
\mathcal M_{10}(\Gamma(1))=\mathbb C G_{10}=\mathbb C G_4G_6.
\]
最后取 $k=12$，每个尖点形式 $f$ 在 $\infty$ 之外都不能为零，且在 $\infty$ 处零点阶数为 $1$.这适用于 $\Delta$，从而
\[
\mathcal S_{12}(\Gamma(1))=\mathbb C\Delta.
\]
现在，$\Delta$ 作为 $\mathbb H$ 上的函数可逆，故映射 $f\mapsto\Delta f$ 对任意 $k\in\mathbb N$ 诱导同构
\[
\mathcal M_k(\Gamma(1))\xrightarrow{\sim}\mathcal S_{k+12}(\Gamma(1)).
\]
另一方面，$\mathcal S_k(\Gamma(1))$ 在 $\mathcal M_k(\Gamma(1))$ 中余维为 $1$（它是 $f\mapsto f(\infty)$ 的核），又由于 $G_k$ 不是尖点形式，对 $k\ge12$ 有
\[
\mathcal M_k(\Gamma(1))=\mathbb C G_k\oplus\mathcal S_k(\Gamma(1)).
\]
\footnote{译注：原文此处的 $k\ge12$ 应结合前文限定在偶数权；奇数权时两空间均为零，不能称余维为 $1$.}
由 $k<12$ 时的计算得到
\[
\dim_{\mathbb C}\mathcal M_k(\Gamma(1))=
\begin{cases}
0,&k\text{ 为奇数},\\
\lfloor k/12\rfloor,&k\equiv2\pmod{12},\\
\lfloor k/12\rfloor+1,&k\equiv0,4,6,8,10\pmod{12}.
\end{cases}
\]
实际上还有更强的结果：
\begin{proposition*}
代数 $\mathcal M(\Gamma(1))$ 是以 $G_4$ 和 $G_6$ 为变量的多项式代数.
\end{proposition*}
\begin{proof}
要证明对 $k\in2\mathbb N$，满足 $n,m\in\mathbb N$ 和 $4n+6m=k$ 的 $G_4^nG_6^m$ 构成 $\mathcal M_k(\Gamma(1))$ 的一组基.$k\le10$ 时已经验证.

对 $k\ge12$ 归纳证明该族生成整个空间.选取 $n,m$ 使 $4n+6m=k$，并取 $f\in\mathcal M_k$.由于 $G_4^nG_6^m$ 在无穷远处不为零，存在 $\lambda$ 使 $f-\lambda G_4^nG_6^m$ 为尖点形式.于是 $f-\lambda G_4^nG_6^m=\Delta g$，其中 $g\in\mathcal M_{k-12}$；由归纳假设得 $f\in\mathbb C[G_4,G_6]$（回顾 $\Delta=60^3G_4^3-27\cdot140^2G_6^2$）.

再证线性无关.取满足 $4n_0+6m_0=k$ 且 $m_0$ 最大的 $(n_0,m_0)$.将线性相关关系除以 $G_4^{n_0}G_6^{m_0}$，得到一个 $\mathbb C[T]$ 中的多项式，以亚纯函数 $G_4^3G_6^{-2}$ 为根；这将迫使该函数为常值，但它并非常值.
\end{proof}
\begin{example*}[{模不变量 $j$}]
由于 $\Delta$ 的权为 $12$，且仅在 $\infty$ 处有一个单零点，函数
\[
j:=1728\frac{g_2^3}{\Delta}=\frac{1728g_2^3}{g_2^3-27g_3^2}
\]
是权为 $0$ 的模函数，在 $\mathbb H$ 上全纯，在 $\infty$ 处有一个一阶极点.特别地，$j$ 下降为全纯函数 $\bar j:Y(1)=\Gamma(1)\backslash\mathbb H\longrightarrow\mathbb C$，它是全纯映射 $\bar j:X(1)\longrightarrow\mathbb P^1(\mathbb C)$ 的限制.
\end{example*}
\begin{proposition*}
$\bar j$ 是 Riemann 曲面的同构.特别地，
\[
\mathcal M_{X(1)}=\mathbb C(\bar j),\qquad
\{\text{级为 }\Gamma(1)\text{、权为 }0\text{ 的模函数}\}=\mathbb C(j).
\]
\end{proposition*}
\begin{proof}
由于 $\bar j$ 只有 $\infty$ 处一个单极点，其次数（或重数）为 $1$，因此它是双全纯映射 $X(1)\xrightarrow{\sim}\mathbb P^1(\mathbb C)$.其余结论随之得出.

［也可利用上面的命题，先看到对每个 $\lambda\in\mathbb C$，函数 $\bar j-\lambda$ 只有一个单零点，再由局部反演定理得出结论.］
\end{proof}
为什么选 $1728$？为了使无穷远处的留数等于 $1$.\footnote{译注：原文称``留数''；这里指以下 $q$ 展开中 $q^{-1}$ 的系数.}可计算出，令 $q=\exp(2i\pi z)$，则 $j(z)=q^{-1}+744+\cdots$.注意 $1728=2^63^3$.


\subsection{关于 $q$ 展开}\label{sec:modular-1-2-5}
令 $q=\exp(2i\pi z)$. 上述模形式或模函数具有一些引人注目的 $q$ 展开；这些展开得到的年代，人们对``特殊函数''的掌握也许比今天更为熟练. 习题课中的一道习题将解释展开式
\begin{align*}
E_k(z):=\frac{1}{2\zeta(k)}G_k(z)=1-\frac{2k}{B_k}\sum_{n=1}^\infty\sigma_{k-1}(n)q^n,
\end{align*}
其中 $B_k$ 是有理数 Bernoulli 数，$\sigma_{k-1}(n)=\sum_{d\mid n}d^{k-1}$. 另一个归于 Jacobi 的经典结果（见下文）是
\begin{align*}
\Delta(z)=(2\pi)^{12}q\prod_{n=1}^\infty(1-q^n)^{24}.
\end{align*}
由此可见，在展开式 $\frac{1}{(2\pi)^{12}}\Delta(z)=\sum_{n=1}^\infty\tau(n)q^n$ 中，函数 $n\mapsto\tau(n)$ 取整数值. Ramanujan 研究了这个函数，并提出以下猜想：
\begin{enumerate}[label=\roman*)]
\item 当 $(m,n)=1$ 时，$\tau(nm)=\tau(n)\tau(m)$；当 $p$ 为素数时，$\tau(p^{r+1})=\tau(p^r)\tau(p)-p^{11}\tau(p^{r-1})$.
\item 当 $p$ 为素数时，$|\tau(p)|\leqslant2\sqrt{p^{11}}$.
\end{enumerate}
后面将证明或解释这些猜想如何得到解决和推广. 最后，Jacobi 的结果还说明，在展开式
\begin{align*}
j(z)=q^{-1}+744+\sum_{n=1}^\infty c(n)q^n
\end{align*}
中，$c(n)$ 都是整数！（这里还要用到 $1/B_2$ 是整数这一事实.）
\begin{remark*}[{关于 Jacobi 公式的一种证明}]
目标是证明函数 $f(z)=q\prod_{n=1}^\infty(1-q^n)^{24}$ 是权为 $12$ 的模形式. 它显然是尖点形式，于是由 $S_{12}(\Gamma(1))=\mathbb C\Delta$ 可知它是 $\Delta$ 的倍数；比较首项便得到因子 $(2\pi)^{12}$. 因为 $f(z)=f(z+1)$，只须证明 $f(-1/z)=z^{12}f(z)$. 值得注意的是，$f$ 的对数导数为
\begin{align*}
\frac{d}{dz}\log(f(z))=2i\pi\left(1-24\sum_{n=1}^\infty\sigma_1(n)q^n\right)=:2i\pi E_2(z).
\end{align*}
若 $\sum_{m,n}(mz+n)^{-2}$ 收敛，就会认出这是权为 $2$ 的归一化 Eisenstein 级数的 $q$ 展开. 沿用计算 $G_k$ 的 $q$ 展开的办法（见习题课），可以发现
\begin{align*}
2\zeta(2)E_2(z)=\sum_{m\in\mathbb Z}\sum_{n\in\mathbb Z_m}\frac{1}{(mz+n)^2},
\end{align*}
其中 $m\ne0$ 时 $\mathbb Z_m=\mathbb Z$，否则 $\mathbb Z_m=\mathbb Z\setminus\{0\}$. 求和顺序对收敛至关重要. 交换求和顺序后仍然收敛，但结果会改变. 可以证明（参见 Serre 的书）
\begin{align*}
\sum_{m\in\mathbb Z}\sum_{n\in\mathbb Z_m}\frac{1}{(mz+n)^2}=-\frac{2i\pi}{z}+\sum_{n\in\mathbb Z}\sum_{m\in\mathbb Z_n}\frac{1}{(mz+n)^2}.
\end{align*}
这意味着 $E_2(-1/z)=z^2E_2(z)+\frac{12z}{2i\pi}$；再由对数导数还原，便得到 $f(1/z)=z^{12}f(z)$.
\footnote{译注：原文末式写作 $f(1/z)$，照录；此处应与前文所要证明的 $f(-1/z)$ 相比较.}
\end{remark*}
\subsection{模形式空间的维数}\label{sec:modular-1-2-6}
回到更一般的同余子群 $\Gamma$. 为计算奇数权模形式空间的维数，需要引入不规则尖点的概念. 它源于以下观察：若函数 $f$ 满足 $f(\left[\begin{smallmatrix}-1&h\\0&-1\end{smallmatrix}\right]z)=-f(z)$，即 $f(z+h)=-f(z)$，则它以 $2h$ 为周期，且在展开式 $f(z)=\sum_{n\in\mathbb Z}a_nq_{2h}^n$（其中 $q_{2h}=\exp(2i\pi z/2h)$）中，所有偶数 $n$ 对应的 $a_n$ 均为零. 特别地，若它在无穷远处全纯，就必在那里为零.
据此，若满足以下条件，就称 $x$ 为 $(\Gamma,k)$ 的\newterm{modular-irregular-cusp}{不规则尖点}{irregular cusp}：
\begin{enumerate}[label=\roman*)]
\item $k$ 为奇数，且 $-1\notin\Gamma$；
\item 写成 $x=\tau\infty$ 时，群 $\tau^{-1}\Gamma_x\tau$ 包含 $\left[\begin{smallmatrix}-1&h_x\\0&-1\end{smallmatrix}\right]$.
\end{enumerate}
若 $x=\tau\infty$ 不规则而 $f$ 是权为 $k$ 的模函数，则 $f[\tau]_k(z+h_x)=-f[\tau]_k(z)$. 写成 $n_\infty=n_\infty^{\mathrm{reg}}+n_\infty^{\mathrm{irr}}$（$k$ 为偶数时，$n_\infty^{\mathrm{reg}}=n_\infty$ 且 $n_\infty^{\mathrm{irr}}=0$）.
\begin{theorem*}
当 $k>2$ 时，有下列维数公式；在 $k$ 为奇数时，假设 $-1\notin\Gamma$：
\begin{align*}
\dim_{\mathbb C}\mathcal M_k(\Gamma)&=(k-1)(g-1)+\left\lfloor\frac k4\right\rfloor n_2+\left\lfloor\frac k3\right\rfloor n_3+\frac k2n_\infty^{\mathrm{reg}}+\left\lfloor\frac k2\right\rfloor n_\infty^{\mathrm{irr}},\\
\dim_{\mathbb C}\mathcal S_k(\Gamma)&=\dim_{\mathbb C}\mathcal M_k(\Gamma)-n_\infty^{\mathrm{reg}}.
\end{align*}
\end{theorem*}
\begin{remark*}
\begin{enumerate}[label=\roman*)]
\item 奇数权时（因而 $-1\notin\Gamma$），有 $n_2=0$. 作为习题！
\item 利用亏格公式，还可写成
\begin{align*}
\dim_{\mathbb C}\mathcal M_k(\Gamma)&=\frac{kd}{12}-g+1+\left(\frac k4-\left\lfloor\frac k4\right\rfloor\right)n_2\\
&\quad+\left(\frac k3-\left\lfloor\frac k3\right\rfloor\right)n_3+\left(\frac k2-\left\lfloor\frac k2\right\rfloor\right)n_\infty^{\mathrm{irr}}.
\end{align*}
\footnote{译注：此式依原文保留. 与前一定理及亏格公式比较，三个带括号的修正项的符号存在疑误.}
因此，若 $12\mid k$，则 $\dim_{\mathbb C}\mathcal M_k(\Gamma)=kd/12-g+1$.
\end{enumerate}
\end{remark*}
\begin{proof}
偶数情形 $k=2k'$，仿照权为 $2$ 时的办法. 固定 $X(\Gamma)$ 上一个亚纯 $k'$ 次微分形式 $\omega_0$（例如取 $\omega_0=\omega_1^{k'}$，其中 $\omega_1$ 是一次微分形式），记 $f$ 为对应的权 $k$ 模函数. 由公式
\begin{align*}
\operatorname{ord}_P(\pi^*(g)f)&=e_P\operatorname{ord}_{\pi(P)}(g)+e_P\operatorname{ord}_{\pi(P)}(\omega_0)+k'(e_P-1),&&P\in\mathbb H,\\
\operatorname{ord}_{x,h_x}(\pi^*(g)f)&=\operatorname{ord}_{\pi(x)}(g)+\operatorname{ord}_{\pi(x)}(\omega_0)+k',&&x\text{ 为尖点},
\end{align*}
可知映射 $g\mapsto\pi^*(g)f$ 诱导同构
\begin{align*}
L(\operatorname{div}(\omega_0)+k'D_\pi)&\xrightarrow{\sim}\mathcal M_k(\Gamma),\\
L\left(\operatorname{div}(\omega_0)+k'D_\pi-\sum_{P\mapsto\infty}[P]\right)&\xrightarrow{\sim}\mathcal S_k(\Gamma),
\end{align*}
其中 $D_\pi=\sum_{P\in X(\Gamma)}(1-1/e_P)[P]$. 从而 $\dim_{\mathbb C}\mathcal M_k(\Gamma)=\ell(k'K+\lfloor k'D_\pi\rfloor)$，且
\begin{align*}
\lfloor k'D_\pi\rfloor=\sum_{P\mapsto i}\left\lfloor\frac k4\right\rfloor[P]+\sum_{P\mapsto j}\left\lfloor\frac k3\right\rfloor[P]+\sum_{P\mapsto\infty}\frac k2[P].
\end{align*}
再利用 $\deg(k'K+\lfloor k'D_\pi\rfloor)>2g-2$，由 Riemann--Roch 公式算出 $\dim\mathcal M_k(\Gamma)$. 为计算 $\mathcal S_k(\Gamma)$ 的维数，只须将 $\operatorname{div}(\omega_0)$ 换成 $\operatorname{div}(\omega_0)-\sum_{P\mapsto\infty}[P]$. 此时只要 $k>2$，确有 $\deg>2g-2$.
奇数情形中，要采用类似论证，必须先知道存在权为 $k$ 的模函数. 暂且承认存在这样的 $f$（条件是 $-1\notin\Gamma$）. 则 $f^2$ 的权为 $2k$，对应于 $X(\Gamma)$ 上一个亚纯 $k$ 次微分形式 $\omega_0$. 由公式
\begin{align*}
\operatorname{ord}_P(\pi^*(g)^2f^2)&=2e_P\operatorname{ord}_{\pi(P)}(g)+e_P\operatorname{ord}_{\pi(P)}(\omega_0)+k(e_P-1),&&P\in\mathbb H,\\
\operatorname{ord}_{x,h_x}(\pi^*(g)^2f^2)&=2\operatorname{ord}_{\pi(x)}(g)+\operatorname{ord}_{\pi(x)}(\omega_0)+k,&&x\text{ 为尖点},
\end{align*}
可见 $g\mapsto\pi^*(g)f$ 诱导同构
\begin{align*}
L\left(\tfrac12\operatorname{div}(\omega_0)+\tfrac k2D_\pi\right)&\xrightarrow{\sim}\mathcal M_k(\Gamma),\\
L\left(\tfrac12\operatorname{div}(\omega_0)+\tfrac k2D_\pi-\tfrac12\sum_{P\mapsto\infty}[P]\right)&\xrightarrow{\sim}\mathcal S_k(\Gamma).
\end{align*}
对 $\mathcal M_k$，须计算下列除子的整数部分：
\begin{align*}
\tfrac12\operatorname{div}(\omega_0)+\tfrac k2D_\pi=\sum_{P\in X(\Gamma)}\left(\tfrac12\operatorname{ord}_P(\omega_0)+\tfrac k2(1-1/e_P)\right)[P].
\end{align*}
对 $P\in Y(\Gamma)$，取 $\widetilde P\in\mathbb H$ 满足 $P=\pi(\widetilde P)$. 等式
\begin{align*}
2\operatorname{ord}_{\widetilde P}(f)=e\operatorname{ord}_P(\omega_0)+k(e-1)
\end{align*}
说明当 $e=1$ 或 $3$ 时，$\operatorname{ord}_P(\omega_0)$ 为偶数. 因 $n_2=0$，故它对所有 $P\in Y(\Gamma)$ 都为偶数，从而
\begin{align*}
\left\lfloor\tfrac12\operatorname{ord}_P(\omega_0)+\tfrac k2(1-1/e_P)\right\rfloor=\tfrac12\operatorname{ord}_P(\omega_0)+\left\lfloor\tfrac k2(1-1/e_P)\right\rfloor.
\end{align*}
对尖点 $P\in X(\Gamma)$，取 $x\in\mathbb P^1(\mathbb Q)$ 使 $\pi(x)=P$. 等式
\begin{align*}
\operatorname{ord}_{x,2h_x}(f)=\operatorname{ord}_{x,h_x}(f^2)=\operatorname{ord}_P(\omega_0)+k
\end{align*}
说明：
\begin{itemize}
\item 若 $x=\tau\infty$ 规则，则 $f[\tau]_k$ 以 $h_x$ 为周期，故 $\operatorname{ord}_{x,2h_x}(f)=2\operatorname{ord}_{x,h_x}(f)$，于是 $\operatorname{ord}_P(\omega_0)$ 为奇数，且
\begin{align*}
\left\lfloor\tfrac12\operatorname{ord}_P(\omega_0)+\tfrac k2(1-1/e_P)\right\rfloor=\tfrac12\operatorname{ord}_P(\omega_0)+\tfrac k2.
\end{align*}
\item 若 $x=\tau\infty$ 不规则，则 $f[\tau]_k$ 以 $2h_x$ 为周期，且 $f[\tau]_k(z+h_x)=-f[\tau]_k(z)$. 故在展开 $f[\tau]_k(z)=\sum_{n\in\mathbb Z}a_nq^n$（$q=\exp(2i\pi z/2h_x)$）中，偶数 $n$ 对应的 $a_n$ 均为零，因而 $\operatorname{ord}_{x,2h_x}(f)$ 为奇数，最终有
\begin{align*}
\left\lfloor\tfrac12\operatorname{ord}_P(\omega_0)+\tfrac k2(1-1/e_P)\right\rfloor=\tfrac12\operatorname{ord}_P(\omega_0)+\tfrac{k-1}{2}.
\end{align*}
\end{itemize}
于是得到
\begin{align*}
\left\lfloor\tfrac12\operatorname{div}(\omega_0)+\tfrac k2D_\pi\right\rfloor
&=\tfrac12\operatorname{div}(\omega_0)+\sum_{P\mapsto j}\left\lfloor\frac k3\right\rfloor[P]\\
&\quad+\sum_{P\mapsto\infty,\mathrm{reg}}\frac k2[P]+\sum_{P\mapsto\infty,\mathrm{irr}}\frac{k-1}{2}[P].
\end{align*}
当 $k\geqslant3$ 时，它的次数大于 $\frac12\deg(\operatorname{div}(\omega_0))=\frac k2(2g-2)>2g-2$，所以可用简化的 Riemann--Roch 公式，得到所述 $\mathcal M_k(\Gamma)$ 的维数.
\footnote{译注：此处不等式照录；当亏格为 $0$ 或 $1$ 时，末一严格不等式并不成立，须另行检查.}
要计算 $\mathcal S_k(\Gamma)$ 的维数，须求除子 $\frac12\operatorname{div}(\omega_0)+\frac k2D_\pi-\frac12\sum_{P\mapsto\infty}[P]$ 的整数部分. 只有尖点处的讨论受影响，细节留作习题.
\end{proof}
\begin{remark*}[{$k=1$ 时会怎样？}]
在权为 $1$ 的模函数存在的前提下，上一证明说明，记 $\omega_0$ 为 $X(\Gamma)$ 上的亚纯微分形式，则
\begin{align*}
\dim_{\mathbb C}\mathcal M_1(\Gamma)=\ell\left(\tfrac12\left(\operatorname{div}(\omega_0)+\sum_{P\mapsto\infty,\mathrm{reg}}[P]\right)\right).
\end{align*}
这个除子看似不是整除子，实际上却是，其次数为 $g-1+n_\infty^{\mathrm{reg}}/2$. 当 $n_\infty^{\mathrm{reg}}>2g-2$ 时，可以得出 $\dim_{\mathbb C}\mathcal M_1(\Gamma)=n_\infty^{\mathrm{reg}}/2$. 否则 Riemann--Roch 公式不能给出确切维数，但给出下界 $\dim_{\mathbb C}\mathcal M_1(\Gamma)\geqslant n_\infty^{\mathrm{reg}}/2$. 无论如何，$\dim_{\mathbb C}\mathcal S_1(\Gamma)=\dim_{\mathbb C}\mathcal M_1(\Gamma)-n_\infty^{\mathrm{reg}}/2$.
\end{remark*}
\begin{remark*}[{权为 $1$ 的模函数的存在性}]
上面奇数权的证明假设了当 $-1\notin\Gamma$ 时存在权为 $k$ 的模函数. 当然只须对权 $1$ 证明. 以下论证使用一个稍后将见到的事实：紧 Riemann 曲面上次数为 $0$ 的``除子模主除子''群 $\operatorname{Pic}^0(X)$ 是可除群.
因此，若 $P\in X(\Gamma)$，则 $\operatorname{Pic}^0(X)$ 中的元素 $K-(2g-2)[P]$ 可被 $2$ 整除. 于是，若 $\omega_0$ 是 $X(\Gamma)$ 上的亚纯微分形式，就另有一个 $\omega$ 满足
\begin{align*}
\operatorname{div}(\omega)=2(\operatorname{div}(\omega_0)-(2g-2)[P]).
\end{align*}
\footnote{译注：此等式依原文保留. 左端次数应为 $2g-2$，右端次数为 $0$，因而一般情形下不相容；后续存在性论证也须据此复核.}
令 $f_2$ 为满足 $f_2(z)dz=\pi^*(\omega)$ 的权为 $2$、级为 $\Gamma$ 的模函数. 则在 $\mathbb H$ 的每一点 $z$，$\operatorname{ord}_z(f_2)$ 为偶数；因 $\mathbb H$ 单连通，存在其上的亚纯函数 $f:\mathbb H\to\mathbb C$ 使 $f^2=f_2$. 此函数未必具有权 $1$ 的自守性，但至少存在特征 $\chi:\Gamma\to\{\pm1\}$，使对所有 $\gamma\in\Gamma$ 都有 $f[\gamma]_1=\chi(\gamma)f$. 特别地，$f$ 对 $\Gamma':=\ker\chi$ 自守；这是 $\Gamma$ 的指数为 $2$ 的子群，且 $f$ 在尖点处也亚纯.
若 $\Gamma'=\Gamma$，便已达到目的. 否则，只须找到在 $\mathbb H$ 及尖点处亚纯的 $f_0$，满足 $\forall\gamma\in\Gamma,\ f_0(\gamma z)=\chi(\gamma)f_0(z)$. 此时 $f_0f$ 就是 $\Gamma$ 的权 $1$ 模函数. 为此注意，$X(\Gamma')$ 在 $X(\Gamma)$ 上的次数为 $2$（这里用到 $-1\notin\Gamma$，它保证 $\Gamma'$ 在 $\mathrm{PSL}_2(\mathbb R)$ 中的像在 $\Gamma$ 的像中指数为 $2$）. 因此存在 $g\in\mathcal M_{X(\Gamma')}\setminus\mathcal M_{X(\Gamma)}$. 函数 $\pi^*g$ 对 $\Gamma'$ 是权 $0$ 模函数，对 $\Gamma$ 则不是. 于是取 $\gamma\in\Gamma\setminus\Gamma'$，函数 $f_0(z):=\pi^*g(z)-\pi^*g(\gamma z)$ 就满足所需条件.
\end{remark*}
\begin{example*}[{$\Gamma(N)$，$N>1$ 的情形}]
习题课中已经见到，此时 $n_2=n_3=0$ 而 $n_\infty=d_N/N$，其中 $d_N$ 是 $X(N)\to X(1)$ 的次数，等于
\begin{align*}
d_N=[\Gamma(1):\Gamma(N)\{\pm1\}]=\begin{cases}\frac12N^3\prod_{p\mid N}(1-p^{-2}),&N\ne2,\\6,&N=2.\end{cases}
\end{align*}
因此亏格为 $g=1+d_N/12-d_N/2N$. 此外，当 $-1\notin\Gamma(N)$（即 $N>2$）时，所有尖点显然都是规则的，因为 $\left[\begin{smallmatrix}-1&n\\0&-1\end{smallmatrix}\right]\notin\Gamma(N)$. 因而当 $N>2$ 时，维数公式成为
\begin{align*}
\dim_{\mathbb C}\mathcal M_k(\Gamma(N))&=\frac{k-1}{12}d_N+\frac{d_N}{2N},&
\dim_{\mathbb C}\mathcal S_k(\Gamma(N))&=\frac{k-1}{12}d_N-\frac{d_N}{2N}.
\end{align*}
\end{example*}

\subsection{Poincaré 级数}\label{sec:modular-1-2-7}
这里介绍构造级为 $\Gamma$ 的模形式的一般方法. 一个朴素想法是从任意函数 $f$ 出发，构造级数 $\varphi(z)=\sum_{\gamma\in\Gamma}f(\gamma z)/j_\gamma(z)^k$. 公式 $j_{\gamma\gamma'}(z)=j_\gamma(\gamma'z)j_{\gamma'}(z)$ 说明，这样的函数确实应对 $\Gamma$ 具有权 $k$ 的自守性. 问题是，这样的级数不可能收敛，例如有无穷多个 $\gamma$ 满足 $j_\gamma(z)=1$. 为绕过这个问题，记
\begin{align*}
\Gamma_\infty^+:=\{\gamma\in\Gamma,j_\gamma=1\}&=\Gamma\cap\left\{\begin{bmatrix}1&n\\0&1\end{bmatrix},\ n\in\mathbb Z\right\}\\
&=\Gamma\cap\begin{bmatrix}1&1\\0&1\end{bmatrix}^{\mathbb Z}.
\end{align*}
这是 $\Gamma_\infty$ 中指数为 $1$ 或 $2$ 的子群. 从一个已经在 $\Gamma_\infty^+$ 下不变的函数 $f$ 出发，则函数 $\varphi(z)=\sum_{\gamma\in\Gamma_\infty^+\backslash\Gamma}f(\gamma z)/j_\gamma(z)^k$ 更有希望绝对收敛，而在收敛时仍具有权 $k$ 的自守性. 这里 $\gamma$ 遍历 $\Gamma$ 关于 $\Gamma_\infty^+$ 的左陪集代表系.
设 $\Gamma_\infty^+=\left[\begin{smallmatrix}1&h\\0&1\end{smallmatrix}\right]^{\mathbb Z}$（注意，$\infty$ 是规则尖点时 $h=h_\infty$，否则 $h=2h_\infty$）. 在 $\Gamma_\infty^+$ 下不变的典型函数有 $f(z)=\exp(2i\pi nz/h)$，$n\in\mathbb Z$.
\begin{proposition*}
权为 $k\geqslant3$、``特征''为 $n\in\mathbb N$ 的\newterm{modular-poincare-series}{Poincaré 级数}{Poincaré series}
\begin{align*}
\varphi_n(z):=\sum_{\gamma\in\Gamma_\infty^+\backslash\Gamma}j_\gamma(z)^{-k}\exp(2i\pi n\gamma z/h)
\end{align*}
定义模形式 $\varphi_n\in\mathcal M_k(\Gamma)$. 此外：
\begin{enumerate}[label=\roman*)]
\item 若 $n>0$，则 $\varphi_n\in\mathcal S_k(\Gamma)$.
\item 当 $n=0$ 时，$\varphi_0$ 在所有尖点 $x\ne\infty$ 处为零，且
\begin{align*}
\varphi_0(\infty)=\begin{cases}[\Gamma_\infty^+\backslash\Gamma_\infty],&k\text{ 为偶数，或 }\Gamma_\infty^+=\Gamma_\infty,\\0,&\text{否则}.\end{cases}
\end{align*}
\end{enumerate}
\end{proposition*}
\begin{proof}
为寻找方便的 $\Gamma_\infty^+\backslash\Gamma$ 代表系，从等式
\begin{align*}
\begin{bmatrix}1&m\\0&1\end{bmatrix}\begin{bmatrix}a&b\\c&d\end{bmatrix}=\begin{bmatrix}a+mc&b+md\\c&d\end{bmatrix}
\end{align*}
出发. 它说明 $\gamma'\in\Gamma_\infty^+\gamma\Rightarrow(a',b')\equiv(a,b)\pmod h$. 反之，若 $\gamma'$ 满足 $(a',b')\equiv(a,b)\pmod h$ 且 $(c',d')=(c,d)$，则行列式条件 $\det\gamma=\det\gamma'=1$ 蕴含 $a-a'=cm$、$b-b'=dm$，其中 $h\mid m$，因而蕴含 $\gamma'\in\Gamma_\infty^+\backslash\Gamma$.
\footnote{译注：原文此处写作商集 $\Gamma_\infty^+\backslash\Gamma$，照录；按此处论证，似应是陪集 $\Gamma_\infty^+\gamma$. 以下仅按第二行选取代表的论证，对一般 $\Gamma$ 亦有待核查.}
最后，对每个 $(c,d)\in\mathbb Z^2$ 选取一个第二行为 $(c\ d)$ 的 $\gamma_{c,d}\in\Gamma$（在能选取时！），便得到 $\Gamma_\infty^+\backslash\Gamma$ 的代表系. 由此
\begin{align*}
\sum_{\Gamma_\infty^+\backslash\Gamma}|j_\gamma(z)^{-k}\exp(2i\pi n\gamma z/h)|\leqslant\sum_{(c,d)\ne(0,0)}\frac{1}{|cz+d|^k}.
\end{align*}
已经见到，$k\geqslant3$ 时右边在一个基本域上一致收敛. 故 Poincaré 级数在各基本域上正规收敛，因而在紧集上正规收敛，故在 $\mathbb H$ 上全纯.
对 $\alpha\in\mathrm{SL}_2(\mathbb Z)$，公式
\begin{align*}
\varphi_n(\alpha z)&=\sum_{\gamma\in\Gamma_\infty^+\backslash\Gamma}j_\gamma(\alpha z)^{-k}q_h^{n\gamma\alpha z}\\
&=j_\alpha(z)^k\sum_{\gamma\in\Gamma_\infty^+\backslash\Gamma}j_{\gamma\alpha}(z)^{-k}q_h^{n\gamma\alpha z}\\
&=j_\alpha(z)^k\sum_{\gamma\in\Gamma_\infty^+\backslash\Gamma\alpha}j_\gamma(z)^{-k}q_h^{n\gamma z}
\end{align*}
在 $\alpha\in\Gamma$ 时说明 $\varphi_n$ 确实是权为 $k$、级为 $\Gamma$ 的自守函数.
\footnote{译注：此处指数记号 $q_h^{n\gamma z}$ 依原文保留；其所表达的是前文的 $\exp(2i\pi n\gamma z/h)$.}
因 $\gamma\notin\Gamma(1)_\infty$ 时，$\Im(z)\to\infty$ 蕴含 $j_\gamma(z)^{-k}\to0$，基本域上的正规收敛还给出：对 $\alpha\in\Gamma(1)$，
\begin{align*}
\lim_{\Im(z)\to\infty}\varphi_n[\alpha]_k(z)=\sum_{\gamma\in\Gamma_\infty^+\backslash(\Gamma\alpha\cap\Gamma(1)_\infty)}\lim_{\Im(z)\to\infty}\bigl(j_\gamma(z)^{-k}\exp(2i\pi n\gamma z/h)\bigr).
\end{align*}
此处的和是有限和，故 $\varphi_n[\alpha]_k$ 在 $\infty$ 处全纯，等价地说，$\varphi_n$ 在尖点 $x=\alpha\infty$ 处全纯. 分两种情形讨论.
若尖点 $x$ 不与 $\infty$ 关于 $\Gamma$ 等价，即 $\infty\notin\Gamma x$，则 $\Gamma\alpha\cap\Gamma(1)_\infty=\varnothing$，所以上式对所有 $n\geqslant0$ 显然为零.
在尖点 $\infty$，可取 $\alpha=1$；此时 $n>0$ 时上述极限为零，而 $n=0$ 时得到命题中的 $\varphi_0(\infty)$ 公式，因为这里当 $\gamma\in\Gamma_\infty\setminus\Gamma_\infty^+$ 时 $j_\gamma(z)=-1$，当 $\gamma\in\Gamma_\infty^+$ 时 $j_\gamma=1$.
\end{proof}
\begin{remark*}[{与 Eisenstein 级数的关系}]
设 $\Gamma=\Gamma(1)$（因而 $k$ 为偶数）且 $n=0$. 与上面的证明一样，可以取 $\Gamma(1)_\infty^+\backslash\Gamma(1)$ 的代表为 $\left[\begin{smallmatrix}a&b\\c&d\end{smallmatrix}\right]$，其中 $(c,d)$ 遍历 $\Gamma(1)$ 中矩阵所有可能的第二行，即所有最大公因数为 $1$ 的 $(c,d)$. 故
\begin{align*}
\varphi_0(z)=\sum_{\gamma\in\Gamma(1)_\infty^+\backslash\Gamma(1)}j_\gamma(z)^{-k}=\sum_{\gcd(c,d)=1}\frac{1}{(cz+d)^k}.
\end{align*}
这与 $G_k$ 尚有差别，但可将 $G_k$ 重排为
\begin{align*}
G_k(z)&=\sum_{(c,d)\ne(0,0)}\frac{1}{(cz+d)^k}=\sum_{n=1}^\infty\sum_{\gcd(c,d)=1}\frac{1}{(cnz+dn)^k}\\
&=\sum_{n=1}^\infty\frac{1}{n^k}\sum_{\gcd(c,d)=1}\frac{1}{(cz+d)^k},
\end{align*}
得到 $G_k(z)=\zeta(k)\varphi_0(z)$，最终 $\varphi_0(z)=2E_k(z)$.
\end{remark*}
\paragraph{应用：$\mathcal M_k(\Gamma)$ 的 Eisenstein 部分.}
设 $\Gamma$ 是 $\Gamma(1)$ 的正规子群，且尖点都是规则的，例如 $N>2$ 时的 $\Gamma(N)$. 对每个 $\alpha\in\Gamma(1)$，函数 $\varphi_0[\alpha]_k$ 属于 $\mathcal M_k(\alpha^{-1}\Gamma\alpha)=\mathcal M_k(\Gamma)$，且除 $x:=\alpha^{-1}\infty$ 外，在其他尖点都为零. 若 $\alpha'\in\Gamma(1)$ 也满足 $x=\alpha'^{-1}\infty$，则存在 $\sigma\in\Gamma(1)_\infty$ 使 $\alpha'=\sigma\alpha$. 由 $j_{\sigma\gamma}(z)=j_\sigma(\gamma z)j_\gamma(z)=j_\sigma j_\gamma(z)$，其中 $j_\sigma=\pm1$（依 $\sigma\in\pm\left[\begin{smallmatrix}1&1\\0&1\end{smallmatrix}\right]^{\mathbb Z}$ 而定），算得
\begin{align*}
\varphi_0[\alpha']_k(z)&=\sum_{\gamma\in\Gamma_\infty^+\backslash\Gamma\alpha'}j_\gamma(z)^{-k}=\sum_{\gamma\in\Gamma_\infty^+\backslash\sigma\Gamma\alpha}j_\gamma(z)^{-k}\\
&=j_\sigma^{-k}\varphi_0[\alpha]_k(z).
\end{align*}
这说明，$k$ 为奇数时 $\varphi_0[\alpha]_k$ 只依赖于 $\alpha$ 在 $\Gamma\backslash\Gamma(1)/\Gamma(1)_\infty^+$ 中的类；$k$ 为偶数时则只依赖于 $\Gamma\backslash\Gamma(1)/\Gamma(1)_\infty$ 中的类，即尖点 $x$. 若再对 $\alpha=\left[\begin{smallmatrix}a&b\\c&d\end{smallmatrix}\right]$ 要求 $d>0$，则对所有 $k$，$\varphi_0[\alpha]_k$ 都只依赖于尖点 $x$，故可定义相应的``Eisenstein 级数''
\begin{align*}
E_k^x:=[\Gamma_\infty:\Gamma_\infty^+]^{-1}\varphi_0[\alpha]_k.
\end{align*}
因 $E_k^x$ 在 $x$ 处非零，在其他尖点处为零，族 $(E_k^x)_{x\in\Gamma\backslash\mathbb P^1(\mathbb Q)}$ 在 $\mathcal M_k(\Gamma)$ 中线性无关. 又已知 $k\geqslant3$ 时 $\dim\mathcal M_k(\Gamma)-\dim\mathcal S_k(\Gamma)=n_\infty$，故此族张成 $\mathcal S_k(\Gamma)$ 的一个补空间，情形与 $\Gamma=\Gamma(1)$ 相同. 称它为 $\mathcal M_k(\Gamma)$ 的``Eisenstein 部分''.
\begin{example*}[{$\Gamma(N)$ 的 Eisenstein 级数}]
当 $\Gamma=\Gamma(N)$、$N>2$ 时，两个尖点 $x=r/t$ 和 $x'=r'/t'$（分子分母互素，$t,t'>0$）等价，当且仅当 $(r,t)\equiv\pm(r',t')[N]$（习题）. 此时算得
\begin{align*}
E_k^x(z)=\sum_{\substack{\gcd(c,d)=1\\(c,d)\equiv(t,-r)[N]}}\frac{1}{(cz+d)^k}.
\end{align*}
于是 $\sum_xE_k^x$ 是 $\Gamma(1)$ 的 Eisenstein 级数 $E_k$.
\footnote{译注：此断言及上式依原文保留；对于奇数 $k$，选取符号与求和的相容性须另核.}
这提示用阶为 $N$ 的行向量 $\bar v\in\mathbb Z/N\mathbb Z\times\mathbb Z/N\mathbb Z$ 为 Eisenstein 级数编指标. 可以定义
\begin{align*}
E_k^{\bar v}(z)&=\sum_{\substack{\gcd(c,d)=1\\(c,d)\equiv\bar v[N]}}\frac{1}{(cz+d)^k},&
G_k^{\bar v}(z)&:=\sum_{(c,d)\equiv\bar v[N]}\frac{1}{(cz+d)^k}.
\end{align*}
与 $\Gamma(1)$ 情形类似的计算给出公式
\begin{align*}
G_k^{\bar v}(z)=\sum_{n\in(\mathbb Z/N\mathbb Z)^\times}\zeta_{n,N}(k)E_k^{n\bar v}(z),\qquad
\zeta_{n,N}(k)=\sum_{\substack{m=1\\m\equiv n[N]}}^\infty\frac{1}{m^k}.
\end{align*}
\end{example*}
\subsection{Petersson 内积}\label{sec:modular-1-2-8}
回顾 $\mathbb H$ 有一个在 $\mathrm{SL}_2(\mathbb R)$ 下不变的测度，称为``Poincaré 测度''，由下式给出：
\begin{align*}
d\mu(z)=\frac{1}{y^2}dx\,dy=\frac{i}{2\Im(z)^2}dz\,d\bar z.
\end{align*}
若 $f:\mathbb H\to\mathbb C$ 连续有界，则对所有 $\alpha\in\mathrm{SL}_2(\mathbb Z)$，$f$ 在基本域 $D=\{z,|\Re(z)|\leqslant1/2,|z|\geqslant1\}$ 上的积分 $\int_Df(\alpha z)d\mu(z)$ 绝对收敛.
现设 $\Gamma$ 为 $\mathrm{SL}_2(\mathbb Z)$ 的同余子群，$f$ 在 $\Gamma$ 下不变（仍假设有界）. 若 $\Gamma(1)=\bigsqcup_{i=1}^r\Gamma\alpha_i$，则和
\begin{align*}
\int_{X(\Gamma)}f\,d\mu:=\sum_{i=1}^r\int_Df(\alpha_i z)d\mu(z)
\end{align*}
不依赖于代表 $\alpha_i$ 的选取. 换言之，若 $D_\Gamma$ 是 $\Gamma$ 在 $\mathbb H$ 中的基本域，则 $\int_{D_\Gamma}f(z)d\mu(z)$ 存在，且
\begin{align*}
\int_{X(\Gamma)}f\,d\mu=\int_{D_\Gamma}f(z)d\mu(z).
\end{align*}
将此用于常函数 $1$，记 $V_\Gamma:=\int_{X(\Gamma)}d\mu$ 为 $X(\Gamma)$ 的体积. 若 $\Gamma'\subset\Gamma$ 为有限指数子群，则 $V_{\Gamma'}=[\Gamma\{\pm1\}:\Gamma'\{\pm1\}]V_\Gamma$（习题）. 因而一般有 $V_\Gamma=d_\Gamma V_{\Gamma(1)}$，并可计算出\footnote{更一般地，可以证明 $\mathbb H\cup\mathbb P^1(\mathbb R)$ 中一个测地三角形的面积等于 $\pi$ 减去其内角之和.} $V_{\Gamma(1)}=\pi/3$.
现取 $f,g\in\mathcal M_k(\Gamma)$. 关系 $\Im(\gamma z)=\Im(z)|j_\gamma(z)|^{-2}$ 表明函数 $h(z):=\Im(z)^kf(z)\overline{g(z)}$ 在 $\Gamma$ 下不变.
\begin{lemma*}
若 $f,g$ 中至少有一个是尖点形式，则函数 $h$ 在 $\mathbb H$ 上有界.
\end{lemma*}
\begin{proof}
只须证明，对 $\alpha\in\Gamma\backslash\Gamma(1)$，$h(\alpha z)$ 在 $\infty$ 的邻域内有界. 若 $h$ 是尖点 $\alpha\infty$ 的``宽度''（即阶），则有展开 $h(z)=\Im(z)^k\sum_{n=0}^\infty a_nq_h^n$，其中 $q_h=\exp(2i\pi z/h)$. 若有一个形式是尖点形式，则 $a_0=0$，且 $h(x+iy)=O(y^k\exp(-cy))$. 又因 $h(\alpha z)$ 沿水平方向周期，故它在 $\infty$ 的邻域内有界.
\footnote{译注：原文此段同时用 $h$ 表示函数与尖点宽度，并把含共轭的乘积写为单个 $q_h$ 级数，照录；严格论证须分别使用 $f[\alpha]_k$ 和 $g[\alpha]_k$ 的展开.}
\end{proof}
\begin{definition*}
对 $f\in\mathcal S_k(\Gamma)$、$g\in\mathcal M_k(\Gamma)$，称
\begin{align*}
\langle f,g\rangle_\Gamma:=\frac{1}{V_\Gamma}\int_{X(\Gamma)}\Im(z)^kf(z)\overline{g(z)}d\mu(z)
\end{align*}
为 $f$ 与 $g$ 的\newterm{modular-petersson-inner-product}{Petersson 内积}{Petersson inner product}.
\end{definition*}
其在 $\mathcal S_k(\Gamma)\times\mathcal S_k(\Gamma)$ 上的限制是正定 Hermite 内积. 按体积归一化后，对 $\Gamma'\subset\Gamma$ 有 $\langle f,g\rangle_{\Gamma'}=\langle f,g\rangle_\Gamma$.
\begin{theorem*}
设 $f\in\mathcal S_k(\Gamma)$ 的展开为 $f(z)=\sum_{n=1}^\infty a_nq_h^n$，其中 $q_h=\exp(2i\pi z/h)$，$h$ 是尖点 $\infty$ 的宽度. 则对每个 Poincaré 级数 $\varphi_n$，有
\begin{align*}
\langle f,\varphi_n\rangle_\Gamma=\begin{cases}0,&n=0,\\\dfrac{h^k(k-2)!}{(4\pi n)^{k-1}}a_n,&n\geqslant1.\end{cases}
\end{align*}
\footnote{译注：此系数依原文保留；按本节内积定义及下文证明，$n\geqslant1$ 的系数应另含因子 $1/V_\Gamma$.}
\end{theorem*}
这个定理最有意思的推论是，所有尖点形式都是 Poincaré 级数的线性组合. 因为任何与所有 $\varphi_n$（$n>0$）正交的尖点形式都为零. 再加上 Eisenstein 级数，得到：
\begin{corollary*}
$\mathcal S_k(\Gamma)$ 由 $\varphi_n$（$n>0$）张成. 此外，若 $\Gamma$ 正规，则 Eisenstein 部分 $\mathcal E_k(\Gamma)$ 与 $\mathcal S_k(\Gamma)$ 正交，且 $\mathcal M_k(\Gamma)=\mathcal S_k(\Gamma)\oplus\mathcal E_k(\Gamma)$.
\end{corollary*}
\begin{proof}[定理的证明]
由定义 $\varphi_n$ 的级数绝对收敛，可写成（$D_\Gamma$ 表示 $\Gamma$ 的一个基本域）
\begin{align*}
\langle f,\varphi_n\rangle_\Gamma
&=\sum_{\gamma\in\Gamma_\infty^+\backslash\Gamma}\frac1{V_\Gamma}\int_{D_\Gamma}\Im(z)^kf(z)\overline{j_\gamma(z)}^{-k}\exp(-2i\pi n\overline{\gamma z}/h)d\mu(z)\\
&=\sum_{\gamma\in\Gamma_\infty^+\backslash\Gamma}\frac1{V_\Gamma}\int_{D_\Gamma}\Im(\gamma z)^kj_\gamma(z)^kf(z)\exp(-2i\pi n\overline{\gamma z}/h)d\mu(z)\\
&=\sum_{\gamma\in\Gamma_\infty^+\backslash\Gamma}\frac1{V_\Gamma}\int_{D_\Gamma}\Im(\gamma z)^kj_{\gamma^{-1}}(\gamma z)^{-k}f(\gamma^{-1}\gamma z)\exp(-2i\pi n\overline{\gamma z}/h)d\mu(z)\\
&=\sum_{\gamma\in\Gamma_\infty^+\backslash\Gamma}\frac1{V_\Gamma}\int_{D_\Gamma}\Im(\gamma z)^kf(\gamma z)\exp(-2i\pi n\overline{\gamma z}/h)d\mu(z)\\
&=\sum_{\gamma\in\Gamma_\infty^+\backslash\Gamma}\frac1{V_\Gamma}\int_{\gamma D_\Gamma}\Im(z)^kf(z)\exp(-2i\pi n\bar z/h)d\mu(z).
\end{align*}
第二行用到 $\Im(z)=\Im(\gamma z)|j_\gamma(z)|^2$，第三行用到 $1=j_{\gamma^{-1}\gamma}(z)=j_{\gamma^{-1}}(\gamma z)j_\gamma(z)$，第四行用到 $f$ 的自守性. 现在注意，$\bigsqcup_{\gamma\in\Gamma_\infty^+\backslash\Gamma}\gamma D_\Gamma$ 是 $\Gamma_\infty^+$ 的一个基本域，依赖于左陪集代表的选取. 可以选取代表，使这个并集是方便的基本域 $D_h=\{z,0\leqslant\Re(z)\leqslant h\}$（中的一个稠密开集）. 因此可写成
\begin{align*}
\langle f,\varphi_n\rangle_\Gamma=\frac1{V_\Gamma}\int_{y=0}^\infty\int_{x=0}^h y^{k-2}f(x+iy)\exp(-2i\pi nx/h)\exp(-2\pi ny/h)dx\,dy.
\end{align*}
代入展开 $f(x+iy)=\sum_{m>0}a_m\exp(2i\pi mx/h)\exp(-2\pi my/h)$，交换积分与求和，得到
\begin{align*}
\langle f,\varphi_n\rangle_\Gamma=\frac h{V_\Gamma}a_n\int_0^\infty y^{k-2}\exp(-4\pi ny/h)dy.
\end{align*}
再作分部积分即得所述公式.
\end{proof}

